\documentclass[11pt]{article}

% Packages for advanced Math, Physics, graphics, and formatting
\usepackage[utf8]{inputenc}
\usepackage{amsmath,amssymb}
\usepackage{booktabs}
\usepackage{enumitem}
\usepackage{geometry}
\usepackage{xcolor}
\usepackage{hyperref}
\usepackage{tikz}
\usepackage{pgfplots}
\usepackage{siunitx}
\geometry{margin=1in}
\hypersetup{
    colorlinks=true,
    linkcolor=blue,
    citecolor=blue,
    urlcolor=blue
}
\pgfplotsset{compat=1.18}

\title{Physics and Mathematics Handbook}
\author{Subhash}
\date{\today}

\begin{document}

% Title page
\begin{titlepage}
    \centering
    {\large \textbf{A Deep Dive into Scientific Writing with LaTeX}}\\[1.5em]
    {\Huge \textbf{Physics and Mathematics Handbook}}\\[2em]
    {\Large \textbf{Author:} Subhash}\\[1em]
    {\large \textbf{Date:} \today}\\[3em]
    \begin{abstract}
        This handbook is designed as a practical LaTeX project that illustrates key concepts in algebra, calculus, physics, diagrams, tables, and references. It is structured to help learners build a five-page document with a title page, table of contents, equations, and illustrative diagrams.
    \end{abstract}
    \vfill
    \textit{A sample learning project for intermediate LaTeX users.}
\end{titlepage}

\tableofcontents
\newpage

% Page 2: Algebra
\section{Algebra}
\subsection{Quadratic Formula}
The quadratic formula solves $ax^2 + bx + c = 0$ with
\begin{equation}
    x = \frac{-b \pm \sqrt{b^2 - 4ac}}{2a}.
\end{equation}
For example, if $a = 1$, $b = 0$, and $c = -4$, then
\begin{equation}
    x = \frac{0 \pm \sqrt{16}}{2} = \pm 2.
\end{equation}

\subsection{Binomial Theorem}
The binomial theorem expands powers of sums:
\begin{equation}
    (x + y)^n = \sum_{k=0}^{n} \binom{n}{k} x^{n-k} y^k.
\end{equation}
For $n=2$:
\begin{equation}
    (x+y)^2 = x^2 + 2xy + y^2.
\end{equation}

\subsection{System of Equations}
A linear system can be written as:
\begin{align}
    2x + 3y &= 5,\\
    4x - y &= 11.
\end{align}
Solving by elimination gives:
\begin{align}
    2(2x + 3y) &= 10,\\
    4x - y &= 11,
\end{align}
so $4x + 6y = 10$ and subtracting yields $7y = -1$, therefore
\begin{equation}
    y = -\frac{1}{7}, \qquad x = \frac{33}{7}.
\end{equation}

\newpage

% Page 3: Calculus
\section{Calculus}
\subsection{Derivatives}
The derivative of a function measures its instantaneous rate of change. For example:
\begin{equation}
    \frac{d}{dx} \left( x^3 - 5x^2 + 4x \right) = 3x^2 - 10x + 4.
\end{equation}

\subsection{Integrals}
The definite integral computes signed area under a curve:
\begin{equation}
    \int_{0}^{2} (3x^2 - 10x + 4) \, dx = \left[ x^3 - 5x^2 + 4x \right]_{0}^{2} = -4.
\end{equation}

\subsection{Limits}
A limit describes behavior as a variable approaches a value:
\begin{equation}
    \lim_{x \to 0} \frac{\sin x}{x} = 1.
\end{equation}

\subsection{Taylor Series}
The Taylor series for $e^x$ around $x=0$ is:
\begin{equation}
    e^x = \sum_{n=0}^{\infty} \frac{x^n}{n!} = 1 + x + \frac{x^2}{2!} + \frac{x^3}{3!} + \cdots.
\end{equation}

\newpage

% Page 4: Physics
\section{Physics}
\subsection{Newton's Laws}
Newton's second law is expressed as
\begin{equation}
    \mathbf{F} = m\mathbf{a}.
\end{equation}

\subsection{Kinetic Energy}
The kinetic energy of a mass $m$ moving at velocity $v$ is
\begin{equation}
    KE = \frac{1}{2} m v^2.
\end{equation}
For $m = \SI{4}{kg}$ and $v = \SI{6}{m/s}$:
\begin{equation}
    KE = \frac{1}{2} (4)(6)^2 = \SI{72}{J}.
\end{equation}

\subsection{Momentum}
Linear momentum is defined by
\begin{equation}
    \mathbf{p} = m \mathbf{v}.
\end{equation}

\subsection{Ohm's Law}
Ohm's law relates voltage, current, and resistance:
\begin{equation}
    V = IR.
\end{equation}
If $V_s = \SI{12.0}{V}$ and $R = \SI{6.0}{\ohm}$, then
\begin{equation}
    I = \frac{V_s}{R} = \SI{2.00}{A}.
\end{equation}

\subsection{Wave Equation}
The one-dimensional wave equation is
\begin{equation}
    \frac{\partial^2 u}{\partial t^2} = c^2 \frac{\partial^2 u}{\partial x^2}.
\end{equation}

\begin{table}[h!]
    \centering
    \caption{Common Physics Quantities}
    \begin{tabular}{@{} lcc @{}}
        \toprule
        Quantity & Symbol & Typical Unit \\
        \midrule
        Mass & $m$ & kg \\
        Velocity & $v$ & m/s \\
        Force & $F$ & N \\
        Energy & $E$ & J \\
        Resistance & $R$ & $\Omega$ \\
        \bottomrule
    \end{tabular}
\end{table}

\newpage

% Page 5: TikZ drawings and conclusion
\section{Diagrams and Visualization}
\subsection{Coordinate System and Parabola}
\begin{figure}[h!]
    \centering
    \begin{tikzpicture}[scale=1]
        \draw[->] (-1,0) -- (4,0) node[right] {$x$};
        \draw[->] (0,-1) -- (0,4) node[above] {$y$};
        \draw[domain=-0.5:2.5,smooth,variable=\x,blue,thick]
            plot ({\x},{0.5*\x^2});
        \node[below left] at (0,0) {$0$};
        \node[below] at (2,0) {$2$};
        \node[left] at (0,2) {$2$};
        \draw[dashed,gray] (2,0) -- (2,2) -- (0,2);
        \node[right] at (2.5,3.2) {$y=\tfrac{1}{2}x^2$};
    \end{tikzpicture}
    \caption{Graph of a simple parabola in a coordinate system.}
\end{figure}

\section{Conclusion}
This handbook demonstrates how LaTeX can produce structured academic content with equations, tables, graphics, and references. Progressing through these examples gives a strong foundation for creating longer technical documents.

\section*{References}
\begin{thebibliography}{9}
    \bibitem{lamport94}
    Leslie Lamport,
    \textit{\LaTeX: A Document Preparation System},
    2nd edition, Addison-Wesley, 1994.

    \bibitem{knuth84}
    Donald E. Knuth,
    \textit{The \TeX book},
    Addison-Wesley, 1984.
\end{thebibliography}

\end{document}